\begin{description}
\item[(T12.1)] Kepler's law:
\[
a = \left(\frac{GM_{\mathrm{s}}}{4\pi^2}P^2\right)^{1/3} \tag*{(1.0)}
\]
Gravitational force provides centripetal acceleration:
\begin{align*}
v_{\mathrm{p}} &= \sqrt{\frac{GM_{\mathrm{s}}}{a}} \tag*{(0.5)}\\
v_{\mathrm{p}} &= \left(\frac{2\pi GM_{\mathrm{s}}}{P}\right)^{1/3}
  \tag*{(1.5)}
\end{align*}

\item[(T12.2)] Momentum conservation:
\[
M_{\mathrm{p}}v_{\mathrm{p}} = M_{\mathrm{s}}v_{\mathrm{s}} \tag*{(1.5)}
\]
Observed quantity is $v_0 = v_{\mathrm{s}}\sin i$. Thus,
\[
M_{\mathrm{p,\,min}} = M_{\mathrm{p}}\sin i
 = \frac{M_{\mathrm{s}}v_{\mathrm{s}}\sin i}{v_{\mathrm{p}}}
 = \frac{M_{\mathrm{s}}v_0}{v_{\mathrm{p}}}
\]
This is a lower limit on $M_{\mathrm{p}}$. \hfill(2.5)

\item[(T12.3)]
\[
bR_{\mathrm{s}} = a\cos i \tag*{(1.0)}
\]
Therefore, for visibility,
$0 \leq b \leq 1 \Rightarrow i \geq \cos^{-1}(R_{\mathrm{s}}/a)$ \hfill(1.0)

\item[(T12.4)] Blackbody $\Rightarrow$ brightness is proportional to area.
Since the observer is far away from the star-planet system, size of
silhouette of planet on stellar disc is independent of $a$.
\[
\Delta = \left(\frac{R_{\mathrm{p}}}{R_{\mathrm{s}}}\right)^2 \tag*{(1.0)}
\]

\item[(T12.5)] Circular orbit $\Rightarrow$ uniform orbital speed
\[
\Rightarrow\ \frac{t}{P} = \frac{a\phi}{2\pi a} = \frac{\phi}{2\pi}
  \tag*{(1.0)}
\]
where $\phi$ is the angle subtended by the planet at the centre of the star
during transit (over time $t$).
% FIGURE NEEDED: (i) stellar disc of radius R_s with transit chord at height
% bR_s, planet positions 1-4, half-chords l_14/2 and l_23/2 subtending
% hypotenuses R_s + R_p and R_s - R_p; (ii) side-on view of the orbit of
% radius a with contact points 1-4 and angles phi_14, phi_23 at the centre
% (2016_theory_sol.pdf p. 23).
\begin{align*}
(l_{23}/2)^2 &= (R_{\mathrm{s}} - R_{\mathrm{p}})^2 - (bR_{\mathrm{s}})^2\\
l_{23} &= 2R_{\mathrm{s}}\sqrt{(1 - R_{\mathrm{p}}/R_{\mathrm{s}})^2 - b^2}
  \tag*{(2.0)}\\
\sin(\phi_{23}/2) &= \frac{l_{23}/2}{a} \tag*{(2.0)}\\
\Rightarrow\ \phi_{23} &= 2\sin^{-1}\left(\frac{l_{23}}{2a}\right)
 = 2\sin^{-1}\left[\frac{R_{\mathrm{s}}}{a}
   \sqrt{\left(1 - \frac{R_{\mathrm{p}}}{R_{\mathrm{s}}}\right)^2 - b^2}
   \right] \tag*{(1.0)}\\
t_{\mathrm{F}} &= \frac{P}{2\pi}\phi_{23}
 = \frac{P}{\pi}\sin^{-1}\left[\frac{R_{\mathrm{s}}}{a}
   \sqrt{\left(1 - \frac{R_{\mathrm{p}}}{R_{\mathrm{s}}}\right)^2 - b^2}
   \right] \tag*{(1.0)}
\end{align*}
Similarly,
\[
t_{\mathrm{T}} = \frac{P}{\pi}\sin^{-1}\left[\frac{R_{\mathrm{s}}}{a}
  \sqrt{\left(1 + \frac{R_{\mathrm{p}}}{R_{\mathrm{s}}}\right)^2 - b^2}
  \right] \tag*{(1.0)}
\]

\item[(T12.6)] Since $R_{\mathrm{s}} \ll a$, use $\sin^{-1}x \approx x$.
\hfill(2.0)
\begin{align*}
t_{\mathrm{T}} &\approx \frac{P}{\pi}\left[\frac{R_{\mathrm{s}}}{a}
  \sqrt{\left(1 + \frac{R_{\mathrm{p}}}{R_{\mathrm{s}}}\right)^2 - b^2}
  \right]\\
t_{\mathrm{F}} &\approx \frac{P}{\pi}\left[\frac{R_{\mathrm{s}}}{a}
  \sqrt{\left(1 - \frac{R_{\mathrm{p}}}{R_{\mathrm{s}}}\right)^2 - b^2}
  \right]
\end{align*}
Dividing, and putting $R_{\mathrm{p}}/R_{\mathrm{s}} = \sqrt{\Delta}$,
\[
\frac{t_{\mathrm{F}}}{t_{\mathrm{T}}}
 = \left[\frac{(1 - \sqrt{\Delta})^2 - b^2}
              {(1 + \sqrt{\Delta})^2 - b^2}\right]^{1/2} \tag*{(1.0)}
\]
\[
\Rightarrow\ b = \left[\frac{(1 - \sqrt{\Delta})^2
  - \left(\dfrac{t_{\mathrm{F}}}{t_{\mathrm{T}}}\right)^2
    (1 + \sqrt{\Delta})^2}
  {1 - \left(\dfrac{t_{\mathrm{F}}}{t_{\mathrm{T}}}\right)^2}\right]^{1/2}
 = \left[1 + \Delta - 2\sqrt{\Delta}\,
   \frac{1 + \left(\dfrac{t_{\mathrm{F}}}{t_{\mathrm{T}}}\right)^2}
        {1 - \left(\dfrac{t_{\mathrm{F}}}{t_{\mathrm{T}}}\right)^2}
   \right]^{1/2} \tag*{(2.0)}
\]

\item[(T12.7)]
\[
t_{\mathrm{T}} = \frac{P}{\pi}\left[\frac{R_{\mathrm{s}}}{a}
  \sqrt{\left(1 + \frac{R_{\mathrm{p}}}{R_{\mathrm{s}}}\right)^2 - b^2}
  \right] \tag*{(1.0)}
\]
Substituting $b$ and $R_{\mathrm{p}}/R_{\mathrm{s}}$,
\begin{align*}
t_{\mathrm{T}} &= \frac{P}{\pi}\left[\frac{R_{\mathrm{s}}}{a}
  \sqrt{(1 + \sqrt{\Delta})^2 - 1 - \Delta + 2\sqrt{\Delta}\,
   \frac{1 + \left(\dfrac{t_{\mathrm{F}}}{t_{\mathrm{T}}}\right)^2}
        {1 - \left(\dfrac{t_{\mathrm{F}}}{t_{\mathrm{T}}}\right)^2}}
  \right]\\
\Rightarrow\ t_{\mathrm{T}} &= \frac{P}{\pi}\frac{R_{\mathrm{s}}}{a}
  \left[\frac{4\sqrt{\Delta}}
        {1 - \left(\dfrac{t_{\mathrm{F}}}{t_{\mathrm{T}}}\right)^2}
  \right]^{1/2}\\
\Rightarrow\ \frac{a}{R_{\mathrm{s}}}
 &= \frac{2P\Delta^{1/4}}
         {\pi(t_{\mathrm{T}}^2 - t_{\mathrm{F}}^2)^{1/2}} \tag*{(2.0)}
\end{align*}

\item[(T12.8)] From part (T12.7)
\[
a = R_{\mathrm{s}}\frac{2P\Delta^{1/4}}
    {\pi(t_{\mathrm{T}}^2 - t_{\mathrm{F}}^2)^{1/2}}
\]
From part (T12.1)
\[
a = \left(\frac{GM_{\mathrm{s}}}{4\pi^2}P^2\right)^{1/3}
\]
Combining,
\[
\frac{GM_{\mathrm{s}}}{4\pi^2}P^2
 = \left[R_{\mathrm{s}}\frac{2P\Delta^{1/4}}
   {\pi(t_{\mathrm{T}}^2 - t_{\mathrm{F}}^2)^{1/2}}\right]^3 \tag*{(3.0)}
\]
\begin{align*}
\Rightarrow\ \rho_{\mathrm{s}}
 &\equiv \frac{M_{\mathrm{s}}}{4\pi R_{\mathrm{s}}^3/3}
 = \frac{3}{4\pi}\,\frac{4\pi^2}{P^2G}\,
   \frac{8P^3\Delta^{3/4}}
        {\pi^3(t_{\mathrm{T}}^2 - t_{\mathrm{F}}^2)^{3/2}} \tag*{(2.0)}
\end{align*}
\[
\Rightarrow\ \boxed{\rho_{\mathrm{s}}
 = \frac{24}{\pi^2 G}\,
   \frac{P(\Delta)^{3/4}}
        {(t_{\mathrm{T}}^2 - t_{\mathrm{F}}^2)^{3/2}}} \tag*{(1.0)}
\]

\item[(T12.9)] From Kepler's third law (with same mass of host star):
\begin{align*}
\frac{a}{a_\oplus} &= \left(\frac{P}{P_\oplus}\right)^{2/3} \tag*{(1.0)}\\
a &= \left(\frac{50.0}{365.242}\right)^{2/3} \times 1\ \mathrm{AU}
 = \boxed{0.266\ \mathrm{AU}} \tag*{(0.5)}\\
&= 0.266 \times 1.496 \times 10^{11}\ \mathrm{m}
 = \boxed{3.97 \times 10^{10}\ \mathrm{m}} \tag*{(0.5)}
\end{align*}

\item[(T12.10)] Edge-on $\Rightarrow b = 0$ \hfill(1.0)
\[
\frac{t_{\mathrm{F}}}{t_{\mathrm{T}}}
 = \left[\frac{(1 - \sqrt{\Delta})^2 - b^2}
              {(1 + \sqrt{\Delta})^2 - b^2}\right]^{1/2}
 = \frac{1 - \sqrt{\Delta}}{1 + \sqrt{\Delta}}
 = \boxed{0.8519} \tag*{(1.0)}
\]

\item[(T12.11)] From parts (T12.1) and (T12.9)
\[
v_{\mathrm{p}} = \sqrt{\frac{GM_\odot}{a}}
 = \sqrt{\frac{6.674 \times 10^{-11} \times 1.989 \times 10^{30}}
              {3.97 \times 10^{10}}}
 = 57.798\ \mathrm{km\,s^{-1}} \tag*{(1.0)}
\]
Assumption of small planet ($M_{\mathrm{p}} \ll M_{\mathrm{s}}$) is valid
because $\Delta$ is very small; less dense planet would make the assumption
stronger!
\[
M_{\mathrm{p}} = \frac{M_{\mathrm{s}}v_0}{v_{\mathrm{p}}}
 = \frac{1.989 \times 10^{30} \times 0.400}
        {5.7798 \times 10^{4} \times 5.972 \times 10^{24}}\,M_\oplus
 = \boxed{2.30\,M_\oplus} \tag*{(2.0)}
\]
From part (T12.4),
\[
R_{\mathrm{p}} = R_{\mathrm{s}}\sqrt{\Delta}
\]
From part (T12.7),
\begin{align*}
\frac{a}{R_{\mathrm{s}}}
 &= \frac{2P\Delta^{1/4}}{\pi(t_{\mathrm{T}}^2 - t_{\mathrm{F}}^2)^{1/2}}
  = \frac{2P\Delta^{1/4}}
         {\pi t_{\mathrm{T}}(1 - (t_{\mathrm{F}}/t_{\mathrm{T}})^2)^{1/2}}\\
\therefore\ R_{\mathrm{s}}
 &= \frac{a\pi t_{\mathrm{T}}
          (1 - (t_{\mathrm{F}}/t_{\mathrm{T}})^2)^{1/2}}{2P\Delta^{1/4}}
\end{align*}
Combining,
\begin{align*}
R_{\mathrm{p}} &= \frac{a\pi t_{\mathrm{T}}
   (1 - (t_{\mathrm{F}}/t_{\mathrm{T}})^2)^{1/2}\Delta^{1/2}}
   {2P\Delta^{1/4}}
 = \frac{a\pi t_{\mathrm{T}}
   (1 - (t_{\mathrm{F}}/t_{\mathrm{T}})^2)^{1/2}\Delta^{1/4}}{2P}
 \tag*{(2.0)}\\
&= \frac{3.97 \times 10^{10} \times \pi \times \frac{1}{24}
   \times (1 - 0.8519^2)^{1/2} \times (0.0064)^{1/4}}
   {2 \times 50.0 \times 6.371 \times 10^{6}}\,R_\oplus\\
&= \boxed{1.21\,R_\oplus} \tag*{(1.0)}
\end{align*}
Mean density
\[
\rho_{\mathrm{p}} = \frac{M_{\mathrm{p}}}{4\pi R_{\mathrm{p}}^3/3}
 = \frac{2.30}{(1.21)^3}\rho_\oplus = 1.3\rho_\oplus \tag*{(1.0)}
\]
Since mean density is higher than that of Earth, the planet is
\fbox{Rocky}. \hfill(1.0)

\item[(T12.12)]
% FIGURE NEEDED: five schematic transit light curves (flux f vs time t),
% each a flat-bottomed dip; transits 1, 3 and 5 are plain dips, while
% transits 2 and 4 show a short upward spike inside the dip (at the same
% phase in both) whose height is almost equal to the maximum dip
% (2016_theory_sol.pdf p. 27).
\begin{itemize}
\item No change in the first: 0.5
\item Spike in second (width and phase of spike is arbitrary): 1.0
\item Height of spike (almost) equal to maximum dip: 0.5
\item No change in third: 0.5
\item Spike again in fourth: 0.5
\item Same phase of spike in second and fourth: 0.5
\item No change in 5th: 0.5
\end{itemize}

\item[(T12.13)]
% FIGURE NEEDED: schematic transit light curve with limb darkening -- dip
% with curved (non-flat) bottom reaching a central minimum
% (2016_theory_sol.pdf p. 27).
Non-flat bottom with central minimum gets 2.0. Curvature of ingress and
egress are tolerated.
\end{description}
